\section*{Sets and Indices}

\begin{itemize}
    \item $R$: set of restaurants, indexed by $r$.
\end{itemize}

\section*{Parameters}

For each restaurant $r \in R$:
\begin{itemize}
    \item $rev_r$: annual revenue in standard operation (from \texttt{Annual\_Revenue}).
    \item $cost_r$: acquisition cost (from \texttt{Cost}).
    \item $h^{\mathrm{std}}_r$: staff hours if operated in Standard mode (from \texttt{Standard\_Staff\_Hours}).
    \item $h^{\mathrm{prem}}_r$: staff hours if operated in Premium mode (from \texttt{Premium\_Staff\_Hours}).
\end{itemize}

Global parameters:
\begin{itemize}
    \item $B$: total investment budget (code/data: \texttt{investment\_budget}).
    \item $\alpha$: premium revenue uplift fraction (code/data: \texttt{premium\_revenue\_uplift}).
    \item $c^{\mathrm{up}}$: premium upgrade cost per premium-operated restaurant (code/data: \texttt{premium\_upgrade\_cost}).
    \item $H$: total staff-hours limit (code/data: \texttt{staff\_hours\_cap}).
\end{itemize}

Named pairwise restriction:
\begin{itemize}
    \item Restaurant \texttt{Restaurant\_D} and Restaurant \texttt{Restaurant\_A} cannot both be acquired.
\end{itemize}

\section*{Decision Variables}

For each restaurant $r \in R$:
\begin{itemize}
    \item $x_r \in \{0,1\}$ (code: \texttt{x[r]}): equals 1 if restaurant $r$ is acquired, 0 otherwise.
    \item $p_r \in \{0,1\}$ (code: \texttt{p[r]}): equals 1 if restaurant $r$ is operated in Premium mode, 0 otherwise.
\end{itemize}

Because of the coupling constraint below, Premium is only allowed when the restaurant is acquired. Thus an acquired restaurant with $p_r=0$ is operated in Standard mode, and an acquired restaurant with $p_r=1$ is operated in Premium mode.

\section*{Objective}

Maximize total annual revenue:
\[
\max \sum_{r \in R} rev_r \bigl(x_r + \alpha p_r\bigr).
\]

This matches the implemented logic: acquiring restaurant $r$ contributes $rev_r$, and selecting Premium adds an incremental uplift of $\alpha \, rev_r$.

\section*{Constraints}

\begin{enumerate}
    \item \textbf{Mode coupling}
    \[
    p_r \le x_r \qquad \forall r \in R.
    \]

    \item \textbf{Combined acquisition and upgrade budget}
    \[
    \sum_{r \in R} \left(cost_r x_r + c^{\mathrm{up}} p_r\right) \le B.
    \]

    \item \textbf{Total staff-hours cap}
    \[
    \sum_{r \in R} \left[h^{\mathrm{std}}_r (x_r - p_r) + h^{\mathrm{prem}}_r p_r\right] \le H.
    \]
    Equivalently, a selected restaurant uses $h^{\mathrm{std}}_r$ hours in Standard mode and $h^{\mathrm{prem}}_r$ hours in Premium mode.

    \item \textbf{Mutual exclusion between Restaurant D and Restaurant A}
    \[
    x_{\texttt{Restaurant\_D}} + x_{\texttt{Restaurant\_A}} \le 1.
    \]

    \item \textbf{Binary restrictions}
    \[
    x_r \in \{0,1\}, \quad p_r \in \{0,1\} \qquad \forall r \in R.
    \]
\end{enumerate}

\section*{Notes on Implementation}

\begin{itemize}
    \item The implemented model is a binary linear optimization model solved with a MIP solver interface (\texttt{GUROBI\_CMD} through PuLP compatibility).
    \item The objective uses only annual revenue and the Premium uplift; it does not subtract acquisition or upgrade costs from the objective. Those costs enter only through the budget constraint.
    \item The business description says each acquired restaurant must be operated either as Standard or Premium. In the implemented model, this is represented implicitly by $x_r \in \{0,1\}$ and $p_r \le x_r$: if $x_r=1$, then $p_r=0$ means Standard and $p_r=1$ means Premium.
    \item Parameter names in the data packet corresponding to the revised Premium logic are exactly \texttt{premium\_revenue\_uplift}, \texttt{premium\_upgrade\_cost}, and \texttt{staff\_hours\_cap}.
\end{itemize}
